\chapter{Cannabis}
\index{Cannabis}

\section*{Definition and Overview}
Cannabis is a plant whose leaves and flowers contain a range of psychoactive compounds, the most important being \textbf{delta-9-tetrahydrocannabinol (THC)}, which causes the feeling of being high, and \textbf{cannabidiol (CBD)}, which is not intoxicating and is being studied for therapeutic effects. Cannabis can be smoked, vaporised, eaten, or taken as an oil, and the products available vary enormously in strength and composition. It is used recreationally and, in some countries, as a prescribed medicine for conditions such as chronic pain, chemotherapy-induced nausea, and some forms of epilepsy. The drug acts on the body's endocannabinoid system and can have both beneficial and harmful effects, depending on how it is used, how often, at what age, and by whom.

\section*{Epidemiology}
Cannabis is one of the most widely used drugs in the world, with an estimated 200 million or more users globally. In many high-income countries, roughly one in ten adults reports use in the past year, with use peaking among people in their late teens and twenties. Men use cannabis more than women, and use is more common among younger people, those with mental health problems, and people who smoke tobacco. Medicinal cannabis use has grown steadily in jurisdictions where it is legal. Rates of cannabis-use disorder, in which use becomes compulsive and harmful, are estimated at around one in ten regular users, rising to about one in six among people who begin using in their teens.

\section*{Aetiology and Pathophysiology}
Cannabis use is a behavioural and public health issue rather than a conventional disease, but understanding it requires understanding how the drug acts. THC mimics natural \textbf{endocannabinoids}, signalling chemicals that help regulate mood, appetite, pain, memory, sleep, and reward. Key points include:
\begin{itemize}
    \item \textbf{Acute effects:} THC binds to cannabinoid receptors in the brain, producing euphoria, relaxation, altered perception, and increased appetite
    \item \textbf{Dependence:} repeated use down-regulates the reward system, so regular users need more to achieve the same effect and feel unwell when they stop
    \item \textbf{Cannabinoid hyperemesis syndrome:} a puzzling but increasingly recognised syndrome of severe, cyclical vomiting and abdominal pain in some long-term users, relieved only by hot showers and by stopping cannabis
    \item \textbf{Effects on the developing brain:} adolescent use is associated with lasting changes in memory, attention, and executive function
    \item \textbf{Risk of mental illness:} use is linked to anxiety and panic and, in vulnerable people, to the onset or worsening of psychotic illness
\end{itemize}
The average potency of modern cannabis has risen greatly compared with earlier decades, increasing both the desired and the unwanted effects of use.

\section*{Clinical Features}
Signs of acute intoxication include:
\begin{itemize}
    \item Feeling relaxed, euphoric, or talkative
    \item Reddened eyes, dry mouth, and increased appetite
    \item Impaired coordination, slowed reaction time, and poor concentration
    \item Anxiety, panic, or paranoia, especially at high doses or in unfamiliar settings
    \item Nausea, dizziness, or light-headedness
\end{itemize}
Regular or heavy use may lead to:
\begin{itemize}
    \item Dependence and withdrawal symptoms, including irritability, poor sleep, anxiety, and cravings on stopping
    \item Poor memory, reduced motivation, and declining performance at school or work
    \item Worsening of anxiety, depression, or psychosis in susceptible people
    \item Chronic bronchitis, cough, and phlegm when the drug is smoked
    \item Recurrent severe vomiting and abdominal pain in cannabinoid hyperemesis syndrome
\end{itemize}

\section*{Differential Diagnosis}
The effects of cannabis may be confused with other conditions, and a detailed history is essential:
\begin{itemize}
    \item Intoxication with or withdrawal from other substances, such as alcohol, stimulants, or sedatives
    \item Anxiety and panic disorders
    \item Bipolar disorder or psychotic illness
    \item Depressive disorders
    \item Migraine or cyclical vomiting syndrome from other causes
    \item Acute gastroenteritis when vomiting dominates the picture
\end{itemize}

\section*{When to Seek Medical Care}
Seek medical advice if cannabis use is causing harm or if you want to stop and are finding it difficult, particularly if:
\begin{itemize}
    \item You feel you cannot control your use despite problems with health, work, or relationships
    \item You experience severe anxiety, panic, or paranoia when using
    \item You have persistent vomiting that you think may be related to cannabis
    \item You are pregnant or planning pregnancy and use cannabis
    \item You have a mental health condition such as schizophrenia or depression
\end{itemize}
\textbf{Contact your local emergency services for:}
\begin{itemize}
    \item Severe agitation, confusion, or hallucinations after cannabis use
    \item A very rapid heart rate with chest pain or fainting
    \item In children: drowsiness, wobbliness, or unresponsiveness after accidentally eating cannabis products
    \item Severe vomiting with signs of dehydration, such as little urine output or dizziness on standing
\end{itemize}

\section*{Diagnosis and Investigations}
\begin{itemize}
    \item \textbf{History:} patterns of use, amount and strength, route of administration, consequences, and a mental health review
    \item \textbf{Examination:} assessment for signs of intoxication or withdrawal
    \item \textbf{Screening questionnaires:} structured tools such as the Cannabis Use Disorders Identification Test (CUDIT)
    \item \textbf{Drug testing:} urine, blood, or saliva tests where needed for medical, occupational, or forensic reasons
    \item \textbf{Mental health assessment:} for depression, anxiety, or psychosis, especially in younger users
    \item \textbf{Review of coexisting problems:} such as respiratory symptoms, pregnancy, or medication interactions
\end{itemize}

\section*{Management}
Management focuses on reducing harm and supporting people who want to cut down or stop:
\begin{itemize}
    \item \textbf{Brief advice and education:} clear, non-judgemental information about the risks, especially for younger users
    \item \textbf{Psychological treatment:} cognitive behavioural therapy and motivational interviewing help people change their use
    \item \textbf{Specialist substance-use support:} counselling and rehabilitation programmes for people with dependence
    \item \textbf{Withdrawal support:} managing cravings, sleep disturbance, and irritability while stopping, often over several weeks
    \item \textbf{Treatment of mental health problems:} for coexisting anxiety, depression, or psychosis
    \item \textbf{Medical cannabis:} in jurisdictions where it is legal, prescribed products with defined doses are used for specific conditions; this is quite different from recreational use
    \item \textbf{Harm reduction:} clear advice against driving, operating machinery, and use during pregnancy or in adolescence
\end{itemize}

\section*{Complications}
\begin{itemize}
    \item Cannabis-use disorder and dependence
    \item Impaired learning, memory, and motivation, especially with adolescent use
    \item Increased risk of anxiety, depression, and psychotic illness in vulnerable people
    \item Cannabinoid hyperemesis syndrome, with recurrent and distressing episodes of vomiting
    \item Respiratory problems from smoking, including chronic cough, bronchitis, and lung irritation
    \item Increased risk of road traffic accidents and workplace injuries while intoxicated
    \item Poorer pregnancy outcomes and possible effects on the developing child
    \item Financial, legal, and relationship problems, and reduced educational achievement
\end{itemize}

\section*{Prognosis and Outcomes}
Many people who use cannabis occasionally experience no serious harm, and a large proportion of young users cut down or stop of their own accord as they take on work and family responsibilities. However, about one in ten regular users develops dependence, and outcomes are worse for those who start young, use very high-strength products, or have pre-existing mental health problems. With psychological support and time, most people who want to stop can do so successfully, and withdrawal symptoms generally settle within a few weeks of stopping.

\section*{Prevention and Self-Care}
\begin{itemize}
    \item Delay or avoid use in adolescence and early adulthood, when the brain is still developing
    \item Avoid cannabis entirely if you have a personal or family history of psychosis or serious mental illness
    \item Do not use cannabis when driving, operating machinery, or during pregnancy or breastfeeding
    \item If you choose to use, prefer lower-potency products and use them less often
    \item Avoid mixing cannabis with alcohol or other sedating drugs
    \item Keep cannabis products securely stored and clearly labelled, out of reach of children
    \item If you want to reduce or stop, seek support early and plan ahead for withdrawal symptoms
\end{itemize}

\section*{Sources and Further Reading}
The following sources provide authoritative, up-to-date information on cannabis, its health effects, and its treatment.
\begin{itemize}
    \item NHS. \textit{Cannabis: health effects and getting help.} https://www.nhs.uk/conditions/
    \item National Institute on Drug Abuse. \textit{Cannabis (marijuana) research report.} https://nida.nih.gov/research-topics/cannabis-marijuana
    \item World Health Organization. \textit{Cannabis: health and social effects.} https://www.who.int/
    \item Centers for Disease Control and Prevention. \textit{Cannabis and health.} https://www.cdc.gov/cannabis/
    \item MedlinePlus. \textit{Cannabis (marijuana).} https://medlineplus.gov/marijuana.html
\end{itemize}
